%==============================================================
%====== FICHIER DE CONFIGURATION DU TRAVAIL DE MATURITÉ =======
%==============================================================
% Vous trouverez les instructions de compilation du fichier
% source principal main.tex dans la documentation se trouvant 
% dans le fichier DocumentationTMlatex.pdf

% Il est indispensable de remplir ce fichier !!!!!!!!!


%---------------------------------------------------------------
%-------------- La page de titre et page de garde --------------
%---------------------------------------------------------------

% Le titre du travail de maturité.
\newcommand{\worktitle}{The axioms of $n$-exact categories}
% Le sous-titre du travail. S'il n'y a pas de soustitre, laissez vide, pour un sous-titre sur plusieurs lignes, mettez un linebreak:
\newcommand{\worksubtitle}{}

% Texte au milieu de la page de titre.
\newcommand{\worktype}{Bachelor's thesis}

% L'auteur/ les auteurs
\newcommand{\theauthor}{Thomas Rosset}
\newcommand{\theadress}{Herman Krags Veg 45}
\newcommand{\lieu}{7050 Trondheim}
%si un seul auteur mettre partout des accolades vides {}
\newcommand{\theauthorB}{}
\newcommand{\theadressB}{}
\newcommand{\lieuB}{}

% La date de soumission du travail.
\newcommand{\workdatemonth}{Spring}
\newcommand{\workdateyear}{2025}

% Le thème du séminaire
\newcommand{\seminar}{}

% Définition du pied de page de titre en fonction de la langue (E nglish, F rancais or D eutsch)
% en décommentant la bonne ligne.
%\newcommand{\titlepagefooter}{\titlepagefooterE}
%\newcommand{\titlepagefooter}{\titlepagefooterD}
\newcommand{\titlepagefooter}{\titlepagefooterF}

\newcommand{\enseignant}{Johanne Haugland}
% Noter éventuellement la(les) branche(s) concernée(s)
\newcommand{\branche}{Abstract Algebra} % commenter cette ligne 
						% s'il y a une branche à préciser
						% et décommenter la suivante
%\newcommand{\branche}{nom de la branche}



%---------------------------------------------------------------
%------- Avant-propos, glossaire, remerciements, index ----------
%---------------------------------------------------------------

% Pour insérer un avant-propos : commentez décommentez Y es ou N o.
% Le fichier à remplir du texte désiré est pagesspeciales/avantpropos.tex
\newcommand{\unavantpropos}{Y}
%\newcommand{\unavantpropos}{N}

% Pour insérer un glossaire : commentez décommentez Y es ou N o.
% Le fichier à remplir du texte désiré est pagesspeciales/glossaire.tex
%\newcommand{\unglossaire}{Y}
\newcommand{\unglossaire}{N}

% Pour insérer des remerciements : commentez décommentez Y es ou N o.
% Le fichier à remplir du texte désiré est pagesspeciales/remerciements.tex
%\newcommand{\desremerciements}{Y}
\newcommand{\desremerciements}{Y}

% Pour insérer un index : commentez décommentez Y es ou N o.
% Il faut alors mettre les clés d'index avec \index{mot} et faire un make index.
%\newcommand{\unindex}{Y}
\newcommand{\unindex}{N}


%---------------------------------------------------------------
%------------ Chapitres et annexes -----------------------------
%---------------------------------------------------------------
% Le nombre de chapitres désirés
\newcommand{\nbchap}{3}
% Le nombre d'annexes désirées
\newcommand{\nbannexes}{}


%---------------------------------------------------------------
%--- Listes des figures, tables, listings, abréviations --------
%---------------------------------------------------------------
% Définition des noms des listes
%\renewcommand{\listfigurename}{Liste des figures} % Nom original = Table des figures
%\renewcommand{\listtablename}{Liste des tables} % Nom original = Liste des tableaux
%\renewcommand{\lstlistlistingname}{Liste des codes sources}  % Nom original = Listings

% Pour insérer une liste des figures : commentez décommentez Y es ou N o.
\newcommand{\unelistefig}{Y}
%\newcommand{\unelistefig}{N}

% Pour insérer une liste des tables : commentez décommentez Y es ou N o.
%\newcommand{\unelistetbl}{Y}
\newcommand{\unelistetbl}{N}

% Pour insérer une liste des listings : commentez décommentez Y es ou N o.
\newcommand{\unelistelst}{N}
%\newcommand{\unelistelst}{N}

% Pour insérer une liste des abréviations/notation : commentez décommentez Y es ou N o.
%\newcommand{\unelisteabreviations}{Y}
\newcommand{\unelisteabreviations}{N}


%---------------------------------------------------------------
%-------------- Objets flottants -------------------------------
%---------------------------------------------------------------
% Définitions des titres des figures, tables, listings...
% English: Figure, Table, Listing (default)
% German: Abbildung, Tabelle, Listing
% French: Figure, Table, Listing
\renewcommand{\figurename}{Figure}
\renewcommand{\tablename}{Table}


%---------------------------------------------------------------
%---------------- Figures --------------------------------------
%---------------------------------------------------------------
% Chemin vers le répertoire des figures. Vous pouvez ajouter d'autres chemins entre accolades
% à l'intérieur des accolades principales.
\graphicspath{{./images/}}

% Insertion de figures ; copier coller dans les chapitres en fonction des besoins
%\tmfigureB{NomFigureSansExtension}{Légende}{fig:votreLabel}	% Taille grande
%\tmfigureN{NomFigureSansExtension}{Légende}{fig:votreLabel}	% Taille normale
%\tmfigureS{NomFigureSansExtension}{Légende}{fig:votreLabel}	% Taille petite
%\tmfigureT{NomFigureSansExtension}{Légende}{fig:votreLabel}	% Taille très petite


%---------------------------------------------------------------
%--------------- Insertion de code : listings ------------------
%---------------------------------------------------------------
% Seul latex est appelé par défaut. On peut charger d'autres languages pour les listings
% selon le modèle de la ligne suivante. Voir le package listings pour les langages disponibles.
%\lstloadlanguages{HTML,PHP,TeX}
% Il faut ensuite placer [...,language=HTML,...] après l'appel \begin{lstlisting} pour définir
% le language à utiliser dans un script particulier


%---------------------------------------------------------------
%--------------- Version provisoire ----------------------------
%---------------------------------------------------------------

% Si vous travaillez avec une version provisoire, cette commande imprime un filigrane
% en haut à gauche de chaque page. Commentez la pour la version finale.
%\reviewtimetoday{\today}{* Draft Version *}


%---------------------------------------------------------------
%---------------- Le texte -------------------------------------
%---------------------------------------------------------------
% Réglage de l'indentation de la première ligne de chaque paragraphe
\parindent=0in		% pas d'indentation
%\parindent=0.1in	% petite indentation
%\parindent=0.2in	% moyenne indentation
%\parindent=0.3in	% grande indentation

% Si vous désirez commenter du code latex sur plusieurs lignes, utilisez la forme :
%\begin{comment} ... \end{comment}


%---------------------------------------------------------------
%----------------- Un index ------------------------------------
%---------------------------------------------------------------

% Il n'est pas nécessaire de réaliser un index. Par défaut, il n'y en a pas. Décommentez si vous en voulez un.
% Pour construire le fichier d'index idx, il doit être compilé avec makeindex en ligne de commande.
%\usepackage{makeidx} \makeindex
